\section{Related Works}

Shopping malls and large retail environments are designed to accommodate
high volumes of customers, often resulting in crowded aisles, dynamic
pedestrian flow, and limited maneuvering space. In such environments,
conventional shopping carts present operational challenges including
limited maneuverability, reduced braking control under heavy loads, and
increased physical strain on users. These challenges are particularly
significant for elderly shoppers and parents managing children, where
divided attention and physical effort increase the risk of minor
collisions and unsafe interactions \cite{goodrich2007, ting2003}.

Studies in retail ergonomics show that pushing heavily loaded carts
significantly increases user exertion and reduces mobility efficiency,
particularly among elderly individuals \cite{wirawan2024}. Additionally,
pedestrian flow research shows that dense crowd conditions require
constant micro-adjustments in direction and speed to avoid collisions,
increasing cognitive load and reaction time demands \cite{molnar1995, ishii2015}.
These conditions can lead to sudden stops, braking instability, and
unintended contact with other shoppers.

From a robotics perspective, indoor environments such as shopping malls
present complex navigation challenges. Unlike structured industrial
settings, retail environments are highly dynamic, with unpredictable
human motion, temporary obstacles, and occlusions \cite{thrun2005}.
Human--robot interaction (HRI) research emphasizes the importance of
adaptive speed control, obstacle prediction, and intuitive user feedback
mechanisms to ensure safety and social acceptance in shared environments
\cite{goodrich2007}.

Furthermore, collision-related incidents involving shopping carts,
particularly heel and ankle impacts, are commonly associated with poor
braking response and momentum under heavy load \cite{smith2006,
vinck2019}. This highlights the need for intelligent control systems
capable of adjusting speed dynamically based on load conditions and
crowd density.

To address these challenges, this project proposes a Smart Autonomous
Shopping Cart Robot designed to aid users through adaptive navigation,
obstacle detection, human-aware motion control, and improved ergonomics.
The system aims to enhance safety, reduce physical strain, and improve
shopping efficiency in crowded retail environments.

Building on this foundation, this paper presents the simulation-based
implementation of the Smart Cart using ROS2 and Gazebo, validating
the proposed control architecture in a physics-accurate virtual
retail environment.